% BHU PYQ -- Manually transcribed & error-corrected
\documentclass[11pt,a4paper]{article}

\usepackage[a4paper,top=2.5cm,bottom=2.5cm,left=2.8cm,right=2.8cm,headheight=14pt]{geometry}
\usepackage{amsmath,amssymb}
\usepackage[T1]{fontenc}
\usepackage[utf8]{inputenc}
\usepackage{lmodern}
\usepackage{microtype}
\usepackage{fancyhdr}
\usepackage{enumitem}
\usepackage{hyperref}
\usepackage{lastpage}
\usepackage{parskip}

\hypersetup{colorlinks=true,urlcolor=black,linkcolor=black,
  pdftitle={CHB-02A: Ancillary Chemistry-I -- BHU PYQ},pdfauthor={Banaras Hindu University}}

\pagestyle{fancy}\fancyhf{}
\renewcommand{\headrulewidth}{0.4pt}\renewcommand{\footrulewidth}{0.4pt}
\fancyhead[L]{\small\textit{Banaras Hindu University}}
\fancyhead[R]{\small\textit{CHB-02A: Ancillary Chemistry-I}}
\fancyfoot[C]{\small Page \thepage\ of \pageref{LastPage}}

\newlist{parts}{enumerate}{2}
\setlist[parts,1]{label=(\alph*),leftmargin=*,itemsep=5pt,topsep=3pt}
\setlist[parts,2]{label=(\roman*),leftmargin=*,itemsep=3pt,topsep=2pt}
\newcommand{\pts}[1]{\hfill\textit{[#1]}}

\begin{document}

\begin{center}
  {\large\bfseries BANARAS HINDU UNIVERSITY}\\[2pt]
  {\normalsize B.Sc. (Hons.) Semester II Examination, 2013-14}\\[6pt]
  \rule{0.8\linewidth}{0.5pt}\\[6pt]
  {\large\bfseries Paper No. CHB-02A: Ancillary Chemistry-I}\\[4pt]
  {\normalsize Time: 3 Hours\hfill Maximum Marks: 70}
\end{center}
\rule{\linewidth}{0.4pt}
\small
\noindent\textit{Instructions: Answer five questions in total, including Question~1 which is compulsory.}
\normalsize
\bigskip

\noindent\textbf{Question 1.} Answer all parts of the following: \pts{5 $\times$ 2 = 10}
\begin{parts}
  \item Indicate the reducing and oxidizing agents in the following reaction:
    \[ \text{CH}_4 + 4\text{Cl}_2 \to \text{CCl}_4 + 4\text{HCl} \]
  \item Write the composition of stainless steel.
  \item What is slag? Write its chemical formula.
  \item Write the names and formulae of two important ores of iron.
  \item Define the term 'isoelectric point' of an amino acid.
\end{parts}

\medskip\rule{\linewidth}{0.2pt}

\noindent\textbf{Question 2.} Differentiate between the following pairs: \pts{7 $\times$ 2 = 14}
\begin{parts}
  \item Distillation and Fractional distillation
  \item TLC (Thin Layer Chromatography) and Paper Chromatography
  \item Primary structure of proteins and secondary structure
  \item Melting point as a purity criterion of a crystalline compound
\end{parts}

\medskip\rule{\linewidth}{0.2pt}

\noindent\textbf{Question 3.} Comment upon the role of metals in our life. How is iron recovered from its ores? Write the chemical reactions involved in the blast furnace. Explain the terms: pig iron, cast iron, and wrought iron. \pts{14}

\medskip\rule{\linewidth}{0.2pt}

\noindent\textbf{Question 4.} Explain the following giving chemical reactions involved: \dots \pts{4 + 4 + 4 + 2 = 14}
\begin{parts}
  \item Preparation and analytical applications of Fehling's solution
  \item Preparation and analytical applications of Benedict's reagent
  \item Determination of glucose in urine using Tollens' reagent
  \item What is the role of Cobalt in the human body?
\end{parts}

\medskip\rule{\linewidth}{0.2pt}

\noindent\textbf{Question 5.} Why are enzymes also called biocatalysts? Describe the different types of enzymes. What is a cofactor? Describe the denaturation of enzymes in very brief. \pts{14}

\medskip\rule{\linewidth}{0.2pt}

\noindent\textbf{Question 6.}
\begin{parts}
  \item Write short notes on: \pts{4 $\times$ 2 = 8}
  \begin{parts}
    \item Mutarotation
    \item Zwitterion
  \end{parts}
  \item Discuss the structure of Nylon-6,6 and Buna-S rubber. Give their monomer units. \pts{6}
\end{parts}

\vfill
\rule{\linewidth}{0.4pt}
\begin{center}\small
  \href{https://bhu-exam.vercel.app/}{Click here to visit our website}\\[4pt]
  \href{https://me-aryan.vercel.app/}{Click here to visit our contributor}\\[4pt]
  \href{https://quashan-me.vercel.app/}{Click here to visit our contributor}
\end{center}

\end{document}
